% Telling the 220 ohm and 10k ohm resistors apart. The resistors are drawn with
% TikZ rather than photographed because the band colors have to stay readable
% when the guide is printed in a classroom.
\definecolor{bodymetal}{HTML}{9FC2E8}
% xcolor's brown is a light tan that reads as gold next to the other bands.
\definecolor{bandbrown}{HTML}{7B4A12}

% \resistorglyph{band 1}{band 2}{band 3}{band 4}{band 5}
% The kit's resistors carry one band on each end bulge and three across the
% middle, all at roughly even spacing, so the geometry here is fixed.
\newcommand{\resistorglyph}[5]{%
\begin{tikzpicture}[x=1cm, y=1cm]
    \draw[line width=1pt, gray!50!black] (-0.5,0.3) -- (2.8,0.3);
    \draw[fill=bodymetal, draw=black!55] (0.24,0.04) rectangle (2.06,0.56);
    \draw[fill=bodymetal, draw=black!55, rounded corners=3pt] (0.02,0) rectangle (0.46,0.6);
    \draw[fill=bodymetal, draw=black!55, rounded corners=3pt] (1.84,0) rectangle (2.28,0.6);
    \fill[#1] (0.16,0.01) rectangle (0.32,0.59);
    \fill[#2] (0.62,0.05) rectangle (0.78,0.55);
    \fill[#3] (1.07,0.05) rectangle (1.23,0.55);
    \fill[#4] (1.52,0.05) rectangle (1.68,0.55);
    \fill[#5] (1.98,0.01) rectangle (2.14,0.59);
\end{tikzpicture}%
}

% \resistorcell{value}{band 1}{band 2}{band 3}{band 4}{band 5}{caption}
\newcommand{\resistorcell}[7]{%
\begin{tabular}{c}
    \textbf{#1} \\[4pt]
    \resistorglyph{#2}{#3}{#4}{#5}{#6} \\[2pt]
    \small #7
\end{tabular}%
}

\begin{tcolorbox}[colback=yellow!10!white,colframe=yellow!50!black]
    \textbf{Which resistor is which?} Both resistors are small and blue with five evenly spaced
    colored bands, and nothing is printed on them. You do not need to figure out which end to read
    from---just count the \textbf{red} bands. The \textbf{220\si{\ohm}} resistor has \textbf{two red
    bands next to each other}, and the \textbf{10\si{\kilo\ohm}} resistor has \textbf{exactly one red
    band}. Swapping the two will not damage anything, but the circuit will not behave the way this
    project describes, so ask an instructor if you are not sure.

    \medskip
    \centering
    \begin{tabular}{c @{\hspace{3em}} c}
        \resistorcell{220\si{\ohm}}{red}{red}{black}{black}{bandbrown}{two red bands} &
        \resistorcell{10\si{\kilo\ohm}}{bandbrown}{black}{black}{red}{bandbrown}{one red band}
    \end{tabular}
\end{tcolorbox}
